\chapter{Tema 2: Señales y Sistemas}

\section{Clasificación de las Señales}

\begin{tcolorbox}[cajaResumen, title={Conceptos Fundamentales: Clasificación}]
Una señal es una función matemática dependiente del tiempo $f(t)$ que describe la variación de una variable física que transporta información codificada.
\begin{itemize}[leftmargin=*]
    \item \textbf{Según el medio físico de transmisión}:
    \begin{itemize}
        \item \textit{Eléctricas}: Voltajes o intensidades transmitidos por cables conductores.
        \item \textit{Electromagnéticas}: Ondas en el espacio libre o guías de onda (radio, satélite).
        \item \textit{Ópticas}: Pulsos luminosos guiados por fibra óptica.
    \end{itemize}
    \item \textbf{Según los valores de la función}:
    \begin{itemize}
        \item \textit{Digitales}: Toman un número finito de valores discretos (ej. niveles binarios $0$ y $V$ voltios).
        \item \textit{Analógicas}: Toman un número infinito de valores continuos (ej. señal telefónica).
    \end{itemize}
    \item \textbf{Según los valores del tiempo}:
    \begin{itemize}
        \item \textit{Continuas}: Definidas para todo instante de tiempo $t$.
        \item \textit{Discretas}: Definidas solo en instantes predeterminados de tiempo.
    \end{itemize}
    \item \textbf{Según su periodicidad}:
    \begin{itemize}
        \item \textit{Periódicas}: Cumplen $f(t) = f(t + T)$, donde $T$ es el período fundamental en segundos.
        \item \textit{No periódicas}: La forma de onda no se repite a lo largo del tiempo.
    \end{itemize}
\end{itemize}
\end{tcolorbox}

\subsection{Caso Parámetrico: Señal Senoidal Pura}
Una señal senoidal pura es la señal periódica fundamental y viene dada por:
\begin{equation}
    f(t) = A \sin(\omega t + \phi) = A \sin(2\pi f t + \phi)
\end{equation}
donde $A$ es la amplitud de pico, $T$ es el período, $f = 1/T$ es la frecuencia en Hz, $\omega = 2\pi/T = 2\pi f$ es la frecuencia angular en rad/s, y $\phi$ es el desfase en radianes.

La \textbf{longitud de onda} ($\lambda$) es la distancia espacial que recorre una onda simple en un período $T$ a la velocidad de propagación $c$:
\begin{equation}
    v_{prop} = \frac{\lambda}{T} \implies \lambda = \frac{c}{f} = c \cdot T
\end{equation}
En el vacío, la velocidad de las ondas electromagnéticas es $c = 3 \cdot 10^8$ m/s.

\section{Dominio del Tiempo y de la Frecuencia}

\begin{tcolorbox}[cajaResumen, title={Representación Temporal vs. Frecuencial}]
\begin{itemize}[leftmargin=*]
    \item \textbf{Dominio del tiempo ($t$)}: Muestra la evolución dinámica de la señal según transcurre el tiempo.
    \item \textbf{Dominio de la frecuencia ($\omega$ o $f$)}: Muestra la amplitud y la fase de las componentes senoidales puras (armónicos) que forman la señal. Un armónico puro $f(t) = V \cos(\omega_0 t)$ genera una sola línea en el dominio de la frecuencia.
\end{itemize}
\end{tcolorbox}

\subsection{Repaso de Números Complejos y Fórmulas de Euler}
\begin{itemize}[leftmargin=*]
    \item \textbf{Unidad imaginaria}: $j = \sqrt{-1}$.
    \item \textbf{Formas binómica y polar}: $z = a + j b = |z| e^{j \phi}$, con $|z| = \sqrt{a^2 + b^2}$ y $\phi = \arctan(b/a)$.
    \item \textbf{Fórmulas de Euler}:
    \begin{equation}
        \cos(x) = \frac{e^{jx} + e^{-jx}}{2}, \quad \sin(x) = \frac{e^{jx} - e^{-jx}}{2j}
    \end{equation}
    \item \textbf{Fase según la posición de un número complejo}:
    \begin{itemize}
        \item Real positivo: $\phi = 0$
        \item Real negativo: $\phi = \pi$
        \item Imaginario puro positivo: $\phi = \frac{\pi}{2}$
        \item Imaginario puro negativo: $\phi = -\frac{\pi}{2}$
    \end{itemize}
\end{itemize}

\section{Serie de Fourier para Señales Periódicas}

Toda función $f(t)$ continua y periódica con período $T$ se descompone en una suma infinita de funciones senoidales puras denominadas armónicos.

\subsection{Forma Trigonométrica}
\begin{equation}
    f(t) = \frac{a_0}{T} + \frac{2}{T} \sum_{n=1}^{\infty} \left( a_n \cos(\omega_n t) + b_n \sin(\omega_n t) \right)
\end{equation}
con $\omega_n = n \omega_1 = \frac{2\pi}{T} n$ y coeficientes:
\begin{equation}
    a_n = \int_{-T/2}^{T/2} f(t) \cos(\omega_n t) \, dt, \quad b_n = \int_{-T/2}^{T/2} f(t) \sin(\omega_n t) \, dt
\end{equation}

Forma compacta en cosenos desfasados:
\begin{equation}
    f(t) = \frac{a_0}{T} + \frac{2}{T} \sum_{n=1}^{\infty} |c_n| \cos(\omega_n t + \theta_n)
\end{equation}
donde $c_n = a_n - j b_n = |c_n| e^{j \theta_n}$, $|c_n| = \sqrt{a_n^2 + b_n^2}$ y $\theta_n = \arctan\left(-\frac{b_n}{a_n}\right)$.

\subsection{Forma Exponencial Compleja}
\begin{equation}
    f(t) = \frac{1}{T} \sum_{n=-\infty}^{\infty} c_n e^{j \omega_n t}, \quad c_n = \int_{-T/2}^{T/2} f(t) e^{-j \omega_n t} \, dt
\end{equation}

\subsection{Espectros Discretos y Teorema de Parseval}
\begin{itemize}[leftmargin=*]
    \item \textbf{Espectro de Amplitud}: Gráfica del módulo $|c_n|$ frente a la frecuencia $\omega$.
    \item \textbf{Espectro de Fase}: Gráfica del argumento $\angle c_n = \theta_n$ frente a la frecuencia $\omega$.
    \item \textbf{Teorema de Parseval}: La potencia media disipada por la señal en una resistencia de $1\,\Omega$ equivale a la suma de las potencias de cada armónico individual:
    \begin{equation}
        \bar{P} = \frac{1}{T} \int_{-T/2}^{T/2} |f(t)|^2 \, dt = \sum_{n=-\infty}^{\infty} \left| \frac{c_n}{T} \right|^2
    \end{equation}
\end{itemize}

\section{Tren de Pulsos Rectangulares Periódico}

Para un tren de pulsos rectangulares de amplitud $V$, anchura $\tau$ y período $T$ centrado en $t=0$:
\begin{equation}
    c_0 = a_0 = V\tau, \quad c_n = V\tau \frac{\sin(\omega_n \tau / 2)}{\omega_n \tau / 2} = V\tau \, \text{sinc}\left(\frac{\omega_n \tau}{2}\right), \quad \theta_n = 0
\end{equation}

\begin{tcolorbox}[cajaResumen, title={Propiedades del Tren de Pulsos}]
\begin{itemize}[leftmargin=*]
    \item Su espectro es discreto con infinitos armónicos a frecuencias múltiplos de $\omega_1 = 2\pi/T$.
    \item La mayor parte de la energía se concentra en las bajas frecuencias bajo la envolvente sinc.
    \item Si el período $T$ disminuye, aumenta la separación entre armónicos ($2\pi/T$).
    \item Si la anchura del pulso $\tau$ disminuye, el contenido en frecuencia se extiende sobre un rango mayor.
\end{itemize}
\end{tcolorbox}

\section{Integral y Transformada de Fourier (Señales No Periódicas)}

Cuando el período tiende a infinito ($T \to \infty$), la separación entre armónicos tiende a cero ($\Delta\omega \to 0$) y el espectro discreto se transforma en un espectro continuo.

\begin{equation}
    F(\omega) = \mathcal{F}\{f(t)\} = \int_{-\infty}^{\infty} f(t) e^{-j\omega t} \, dt
\end{equation}
\begin{equation}
    f(t) = \mathcal{F}^{-1}\{F(\omega)\} = \frac{1}{2\pi} \int_{-\infty}^{\infty} F(\omega) e^{j\omega t} \, d\omega
\end{equation}

\subsection{Pulso Rectangular Aislado}
Para un pulso rectangular de amplitud $V$ y duración $\tau$ centrado en $t=0$:
\begin{equation}
    F(\omega) = V \tau \frac{\sin(\omega \tau / 2)}{\omega \tau / 2}
\end{equation}

\section{Ancho de Banda de una Señal}

El ancho de banda es la medida de la dispersión frecuencial donde se concentra la mayor parte de la potencia de la señal.

\begin{itemize}[leftmargin=*]
    \item \textbf{Criterio 1 (Primer corte por cero)}: Frecuencia dada por el primer cero del espectro. Para un pulso de anchura $\tau$:
    \begin{equation}
        \omega_{AB} = \frac{2\pi}{\tau} \text{ rad/s} \quad \iff \quad f_{AB} = \frac{1}{\tau} \text{ Hz}
    \end{equation}
    \item \textbf{Criterio 2 (Banda limitada)}: Frecuencia máxima a partir de la cual el espectro es estrictamente nulo.
    \item \textbf{Criterio 3 (Frecuencias desplazadas)}: Diferencia entre la frecuencia máxima y mínima ($\Delta f = f_{max} - f_{min}$).
\end{itemize}

\section{Sistemas Lineales Estacionarios (LTI) y Respuesta en Frecuencia}

\subsection{Definición de Distorsión y Tipos de Sistemas}
La distorsión es la alteración en la forma de onda de una señal al atravesar un sistema.
\begin{itemize}[leftmargin=*]
    \item \textbf{Sistemas Lineales}: No generan nuevas componentes de frecuencia; solo modifican la amplitud y la fase de las componentes de entrada.
    \item \textbf{Sistemas No Lineales}: Introducen nuevas componentes frecuenciales en la salida.
\end{itemize}

\subsection{Función de Transferencia}
Un sistema LTI se caracteriza unívocamente por su función de transferencia $H(\omega)$, que relaciona la transformada de Fourier de la salida $G(\omega)$ con la de la entrada $F(\omega)$:
\begin{equation}
    H(\omega) = \frac{G(\omega)}{F(\omega)} \implies G(\omega) = H(\omega) F(\omega)
\end{equation}
Ejemplo para una red eléctrica RC serie (salida en el condensador):
\begin{equation}
    H(\omega) = \frac{1}{1 + j\omega RC}
\end{equation}

\subsection{Teorema de la Respuesta Lineal}
\begin{itemize}[leftmargin=*]
    \item \textbf{Entrada senoidal trigonométrica}: Si $f(t) = A \cos(\omega_1 t)$, la salida es:
    \begin{equation}
        g(t) = A |H(\omega_1)| \cos(\omega_1 t + \theta(\omega_1))
    \end{equation}
    \item \textbf{Entrada exponencial compleja}: Si $f(t) = A e^{j \omega_0 t}$, la salida es:
    \begin{equation}
        g(t) = H(\omega_0) A e^{j \omega_0 t} = |H(\omega_0)| A e^{j(\omega_0 t + \theta(\omega_0))}
    \end{equation}
\end{itemize}

\subsection{Ancho de Banda del Sistema}
El ancho de banda $B_s$ es la frecuencia superior que el sistema deja pasar:
\begin{itemize}[leftmargin=*]
    \item \textbf{Criterio de -3 dB}: Frecuencia $B_s$ donde el módulo decae a $|H(2\pi B_s)| = 0.707 = 1/\sqrt{2}$ de su valor máximo (reducción a la mitad de la potencia).
    \item \textbf{Criterio de primer paso por cero}: Frecuencia en la cual el módulo $|H(\omega)|$ se anula por primera vez.
\end{itemize}

\subsection{Estudio Cualitativo de la Respuesta en Frecuencia}
Dada una señal de ancho de banda $B$ que entra a un sistema de ancho de banda $B_s$:
\begin{itemize}[leftmargin=*]
    \item \textbf{Caso $B > B_s$}: La salida está fuertemente distorsionada. La forma final la impone el sistema y su ancho de banda efectivo pasa a ser $B_s$.
    \item \textbf{Caso $B \sim B_s$}: Distorsión moderada; la salida conserva cierta semejanza con la entrada.
    \item \textbf{Caso $B < B_s$}: Transmisión sin distorsión; la salida es una réplica fiel de la entrada.
\end{itemize}

\section{Perturbaciones en la Transmisión de Señales}

\begin{tcolorbox}[cajaResumen, title={Tipos de Perturbaciones}]
\begin{itemize}[leftmargin=*]
    \item \textbf{Atenuación}: Pérdida inevitable de potencia por la distancia. Reduce la amplitud pero no altera la forma de onda. Se mide en dB:
    \begin{equation}
        N = 10 \log_{10}\left(\frac{P_{enviada}}{P_{recibida}}\right) \text{ dB}
    \end{equation}
    \item \textbf{Distorsión de Atenuación}: Producida cuando el medio atenúa de forma desigual a las diferentes frecuencias. Se corrige mediante ecualizadores de atenuación.
    \item \textbf{Distorsión de Retardo (o Fase)}: Ocurre cuando las componentes de distinta frecuencia se propagan a velocidades diferentes, llegando al receptor en instantes distintos. Se corrige mediante ecualizadores de retardo.
    \item \textbf{Ruido}: Procesos aleatorios aditivos.
    \begin{itemize}
        \item \textit{Térmico o Blanco}: Inevitable, origen agitación térmica. Potencia uniforme en todo el espectro.
        \item \textit{Diafonía (Crosstalk)}: Acoplamiento electromagnético indeseado entre líneas o cables próximos.
        \item \textit{Impulsivo}: Picos de corta duración y elevada amplitud (conmutaciones, tormentas). Muy destructivo en transmisión de datos digitales.
    \end{itemize}
\end{itemize}
\end{tcolorbox}

\subsection{Condiciones para un Medio de Transmisión Ideal}
Para lograr una transmisión totalmente libre de distorsión, la función de transferencia del medio debe cumplir:
\begin{equation}
    |H(\omega)| = K \text{ (constante)}, \quad \theta(\omega) = -\omega t_d \text{ (fase lineal con pendiente negativa)}
\end{equation}
Si $K = 1$, el medio es ideal puro (la señal ni se atenúa ni se distorsiona, solo experimenta un retardo $t_d$).

\section{Formulario Completo del Tema 2}

\begin{tcolorbox}[cajaFormulario, title={Formulario General: Señales y Series de Fourier}]
\textbf{1. Parámetros de Señales Senoidales y Ondas}:
\[
f = \frac{1}{T}, \quad \omega = 2\pi f = \frac{2\pi}{T}, \quad \lambda = \frac{c}{f} = c \cdot T
\]

\textbf{2. Serie Compleja de Fourier (Señales Periódicas)}:
\[
f(t) = \frac{1}{T} \sum_{n=-\infty}^{\infty} c_n e^{j n \omega_1 t}, \quad c_n = \int_{-T/2}^{T/2} f(t) e^{-j n \omega_1 t} \, dt
\]

\textbf{3. Relación entre Coeficientes Trigonométricos y Complejos}:
\[
c_n = a_n - j b_n, \quad |c_n| = \sqrt{a_n^2 + b_n^2}, \quad \theta_n = \arctan\left(-\frac{b_n}{a_n}\right)
\]

\textbf{4. Potencia Media (Teorema de Parseval para Señales Periódicas)}:
\[
\bar{P} = \frac{1}{T} \int_{-T/2}^{T/2} |f(t)|^2 \, dt = \sum_{n=-\infty}^{\infty} \left|\frac{c_n}{T}\right|^2
\]

\textbf{5. Tren de Pulsos Rectangulares Periódico}:
\[
c_0 = V\tau, \quad c_n = V\tau \frac{\sin(\omega_n \tau / 2)}{\omega_n \tau / 2}, \quad \omega_{AB} = \frac{2\pi}{\tau} \text{ rad/s}
\]
\end{tcolorbox}

\begin{tcolorbox}[cajaFormulario, title={Formulario General: Transformada de Fourier y Sistemas LTI}]
\textbf{1. Par de Transformadas de Fourier}:
\[
F(\omega) = \int_{-\infty}^{\infty} f(t) e^{-j\omega t} \, dt, \quad f(t) = \frac{1}{2\pi} \int_{-\infty}^{\infty} F(\omega) e^{j\omega t} \, d\omega
\]

\textbf{2. Propiedades de la Transformada de Fourier}:
\begin{itemize}[leftmargin=*]
    \item Linealidad: $\mathcal{F}\{a f_1(t) + b f_2(t)\} = a F_1(\omega) + b F_2(\omega)$
    \item Componente de continua: $F(0) = \int_{-\infty}^{\infty} f(t) \, dt$
    \item Escalado temporal: $\mathcal{F}\{f(at)\} = \frac{1}{|a|} F\left(\frac{\omega}{a}\right)$
    \item Desplazamiento temporal: $\mathcal{F}\{f(t - t_0)\} = F(\omega) e^{-j \omega t_0}$
    \item Modulación frecuencial: $\mathcal{F}\{f(t) e^{j \omega_0 t}\} = F(\omega - \omega_0)$
    \item Modulación cosenoidal: $\mathcal{F}\{f(t) \cos(\omega_0 t)\} = \frac{F(\omega - \omega_0) + F(\omega + \omega_0)}{2}$
    \item Derivación: $\mathcal{F}\left\{\frac{df(t)}{dt}\right\} = j \omega F(\omega)$
    \item Integración: $\mathcal{F}\left\{\int_{-\infty}^{t} f(\tau) d\tau\right\} = \frac{F(\omega)}{j\omega}$
\end{itemize}

\textbf{3. Respuesta de Sistemas LTI}:
\[
G(\omega) = H(\omega) F(\omega), \quad H(\omega) = |H(\omega)| e^{j\theta(\omega)}
\]

\textbf{4. Medidas de Potencia y Atenuación en Decibelios}:
\[
\text{Relación de Potencias (dB)} = 10 \log_{10}\left(\frac{P_A}{P_B}\right)
\]
\[
\text{Atenuación (dB)} = 10 \log_{10}\left(\frac{P_{enviada}}{P_{recibida}}\right)
\]
\end{tcolorbox}